\section{等价关系与商集}

\begin{enumerate}
\def\labelenumi{\arabic{enumi}.}
\tightlist
\item
  令 \((\mathbf{A}_t)_{i \in \mathbf{I}}\) 是集合 \(\mathbf{E}\) 的一个划分。关系 \(\mathbf{R}\{x, y\} :\) “存在 \(i \in \mathbf{I}\) 使得 \(x \in \mathbf{A}_t\) 并且 \(y \in \mathbf{A}_t\)''（两个一般元素 \(x, y\) （ \(\mathbf{E}\) 的）之间的）满足以下条件：
\end{enumerate}

\begin{enumerate}
\def\labelenumi{(\alph{enumi})}
\item
  \(\mathbf{R}\{x,x\}\) 是恒等式（自反性） \(\mathbf{R}\) ).
\item
  \(\mathbf{R}\{x,y\}\) 并且 \(\mathbf{R}\{y,x\}\) 是等价的（对称性 \(\mathbf{R}\) ).
\item
  关系“R\{x,y\} 和 R\{y,z\}'' 蕴含 R\{x,z\} （R 的传递性）。
\end{enumerate}

如果 C 表示由关系 R 定义的 \(E \times E\) 的子集，则条件 (a)、(b)、(c) 分别等价于以下条件： \((a') \Delta \subset C; (b') \overline{C} = C; (c') C \circ C \subset C.\) 由 \((a')\) 并且 \((c')\) 可得 \(C \circ C = C.\)

\begin{enumerate}
\def\labelenumi{\arabic{enumi}.}
\setcounter{enumi}{1}
\tightlist
\item
  反之，设 \(R\{x,y\}\) 是自反、对称且传递的关系，并设C为其在 \(E \times E\) 。则 E 在该映射下的像 F \(x \to C(x)\) 将E映入 \(\mathfrak{P}(E)\) 是 E 的一个划分，且关系“存在一个子集 \(X \in F\) 使得 \(x \in X\) 并且 \(y \in X\) ”等价于 \(R\{x, y\}\) .
\end{enumerate}

满足条件(a)、(b)和(c)的每一个关系都称为E上的等价关系。该划分 \(\mathfrak{F}\) 它所定义的，作为 \(\mathfrak{P}(\mathbf{E})\) 的子集，称为 E 关于关系 R 的商集，并记为 E/R；其元素称为关于 R 的等价类。映射 \(x \to \mathrm{C}(x)\) 从 E 到 E/R，将每个元素 \(x\) 映到包含 \(x\) 的等价类，称为 E 到 E/R 的典范映射。

相等关系 x = y 是一个等价关系。E 到相应商集的典范映射正是 \(x \to \{x\}\) ，并且是双射。

若 R 是一个等价关系，记号“ \(x \equiv y \pmod{R}\) “”有时用作R的同义词 \(\{x, y\}\) ；它被读作“ \(x\) 等价于 \(y\) 模R”。

\begin{enumerate}
\def\labelenumi{\arabic{enumi}.}
\setcounter{enumi}{2}
\tightlist
\item
  在乘积集合 \(E \times F\) 上，关系“ \(pr_{1}z = pr_{1}z'\) ”在z与 \(z'\) 之间是一个等价关系R，且商集 \((\mathbf{E} \times \mathbf{F}) / \mathbf{R}\) 可以与E建立一一对应关系（这正是“商集”名称的由来）。
\end{enumerate}

更一般地，设 \(f\) 是集合 \(\mathbf{E}\) 到一个集合 \(\mathbf{F}\) 中的一个映射。则关系“ \(f(x) = f(y)\) '' 是 \(\mathbf{E}\) 上的一个等价关系。若我们将此关系记为 \(\mathbf{R}\) ，映射 \(z \to \overline{f}^1(z)\) （其中 \(\overline{f}^1(z)\) 被视为 \(\mathbf{E}/\mathbf{R}\) 的元素）是 \(f(\mathbf{E})\) 到 \(\mathbf{E}/\mathbf{R}\) 上的一个双射.

由此得出，\(f\) 可视为以下三个映射按给定顺序的复合：

\begin{enumerate}
\def\labelenumi{(\arabic{enumi})}
\item
  子集 \(f(\mathbf{E})\) （ \(\mathbf{F}\) 的）到集合 \(\mathbf{F}\) 中的典范映射 ;
\item
  E/R 到 \(f(\mathbf{E})\) 的双射映射，其逆映射刚刚已被定义；
\item
  E 到 E/R 的典范映射。
\end{enumerate}

这种映射的分解称为其典范分解或典范因子分解。

\begin{enumerate}
\def\labelenumi{\arabic{enumi}.}
\setcounter{enumi}{3}
\item
  集合 E 上的每个等价关系 R 都可以通过前一子节中的映射来定义；因为，若 C 是 R 的图，则关系“C(x) = C(y)”等价于 R\{x, y\}.
\item
  设 R 是集合 E 上的一个等价关系，且 A 是 E 的一个子集。则 A 中两个一般元素 x、y 之间的关系 \(R\{x, y\}\) 是 A 上的一个等价关系，称为由 R 在 A 上诱导的关系，并记为 \(R_{A}\) 。设 f 为 E 到 E/R 上的典范映射，g 为 A 到 \(A/R_{A}\) 上的典范映射。通过使 E/R 的一个元素与 \(A/R_{A}\) 的一个元素相对应——当它们分别是 E 的同一元素在 f 和 g 之下的像时——我们便在 A 在 f 之下的像 \(f(A)\) 与商集 \(A/R_{A}\) 之间建立了一个一一对应。如果 \(\varphi\) 表示 A 到 E 中的典范映射，这一对应关系由映射 \(z \to f(\varphi(\overline{g}^1(z)))\) 及其逆映射实现，二者也都称为典范的。
\item
  子集 \(\mathbf{A}\) （ \(\mathbf{E}\) 的）称为关于等价关系 \(\mathbf{R}\) 是饱和的，如果对每个 \(x \in \mathbf{A}\) ，元素 \(x\) 关于 \(\mathbf{R}\) 的等价类包含于 \(\mathbf{A}\) 。换句话说，关于 \(\mathbf{R}\) 的饱和集就是关于 \(\mathbf{R}\) 的等价类的并集。如果 \(f\) 是 \(\mathbf{E}\) 到 \(\mathbf{E} / \mathbf{R}\) 上的典范映射，那么一个集合是饱和的，当它具有形式 \(\overline{f}^1(\mathbf{X})\) ，其中 \(\mathbf{X} \subset \mathbf{E} / \mathbf{R}\) .
\end{enumerate}

设A是E的一个子集。包含A的饱和集的交集是 \(\overline{f}^{1}(f(A))\) 。该集合也可定义为A中元素的等价类的并集，并称为A（关于R）的饱和集。

\begin{enumerate}
\def\labelenumi{\arabic{enumi}.}
\setcounter{enumi}{6}
\tightlist
\item
  令 \(\mathrm{P}\{x, y, z\}\) 是一个关系，其中出现了一个一般元素 \(x\) （属于 \(\mathbf{E}\) ）。那么 \(\mathrm{P}\) 称为（在 \(x\) 中）与等价关系 \(\mathbf{R}\) 相容的，如果关系“ \(\mathrm{P}\{x, y, z\}\) 并且 \(x \equiv x' (\text{mod } \mathbf{R})\) '' 蕴含 \(\mathrm{P}\{x', y, z\}\) .
\end{enumerate}

令 \(f\) 为 \(\mathbf{E}\) 到 \(\mathbf{E} / \mathbf{R}\) 上的典范映射，并且令 \(t\) 为 \(\mathbf{E} / \mathbf{R}\) 的一个一般元素。关系“存在 \(x \in \overline{f}^1(t)\) 使得 \(\mathrm{P}\{x, y, z\}\)'' 则等价于“对所有 \(x \in \overline{f}^1(t)\) , \(\mathrm{P}\{x, y, z\}\)''；后者是 \(t, y, z\) 之间的一个关系，称为由 \(\mathrm{P}\) 在过渡到商（关于 \(x\) ）时诱导的。如果我们把它记为 \(\mathrm{P}'\{t, y, z\}\) ，那么 \(\mathrm{P}\{x, y, z\}\) 等价于 \(\mathrm{P}'\{f(x), y, z\}\) .

对于涉及任意多个参数的关系，以及关系在其多个参数中与 R 相容的情形，也有类似的定义。例如，若 A 是 E 的子集，则说关系“ \(x \in A\) ”是相容的（在 \(x\) ）与 R 等价于说 A 关于 R 是饱和的。如果 \(\varphi\) 是从 E 到集合 F 的映射，那么说函数关系“ \(y = \varphi(x)\) ”是相容的（在 \(x\) ）与 R 成立，就是说 \(\varphi\) 在关于 R 的每个等价类上是常值的。通过过渡到商集，R 因此在 E/R 的 \(y\) 以及一个一般元素 \(t\) 之间诱导出一个关系；这个关系在 \(y\) 中是函数性的，从而确定了一个映射 \(\varphi'\) 从E/R到F，满足恒等式 \(\varphi(x) = \varphi'(f(x))\) .

\begin{enumerate}
\def\labelenumi{\arabic{enumi}.}
\setcounter{enumi}{7}
\item
  设 R 是集合 E 上的一个等价关系，设 S 是集合 F 上的一个等价关系，并设 f 是从 E 到 F 的一个映射。若关系 \(x \equiv x' \pmod{R}\) 蕴含 \(f(x) \equiv f(x') \pmod{S}\) . 如果 g 是 F 到 F/S 的典范映射，那么复合函数 \(g \circ f\) 在关于 R 的等价类 z 的所有元素上取值相同；如果我们用 \(h(z)\) 表示这个公共值，那么 h 是 E/R 到 F/S 的一个映射，并称 h 是在取商时由 f 诱导的。
\item
  设R是E上的一个等价关系，并设S是E/R上的一个等价关系。若f是E到E/R的典范映射，则 \(f(x) \equiv f(y) \pmod{S}\) '' 是 E 上的一个等价关系 T。因此，关于 T 的一个等价类是 E 中关于 R 的等价类中那些关于 S 彼此等价的等价类在 E 中的并集；而关系 `` \(x \equiv y \pmod{R}\) '' 蕴含 `` \(x \equiv y \pmod{T}\) ''。如果 \(g\) 并且 \(\varphi\) 分别是E/R到(E/R)/S以及E到E/T的典范映射，我们通过将(E/R)/S中的一个元素与E/T中的一个元素对应起来（当它们分别是E中同一元素在 \(g \circ f\) 与 \(\varphi\) 下的像时），就得到一个一一对应。
\end{enumerate}

反之，设R和T是E上的两个等价关系，使得 \(x \equiv y \pmod{R}\) '' 蕴含 `` \(x \equiv y \pmod{T}\) ''. 则T与R相容（按第7条的意义），在 \(x\) 以及 \(y\) 两者中，通过过渡到商集E/R（关于 \(x\) 以及 \(y\) ），T在E/R上诱导出一个等价关系S。若 \(f\) 再次表示E到E/R的典范映射，关系“ \(f(x) = f(y) \pmod{S}\) ”等价于“ \(x \equiv y \pmod{T}\) ''。等价关系 S 称为 T 关于 R 的商，并记为 T/R。由前一段可知，在 (E/R)/(T/R) 与 E/T 之间存在一一对应（典范对应）。

\begin{enumerate}
\def\labelenumi{\arabic{enumi}.}
\setcounter{enumi}{9}
\tightlist
\item
  现在设E、F为任意两个集合，它们可以相同也可以不同。设R\{x, y\} 是 E 上的一个等价关系，并令 S\{z, t\} 是 F 上的一个等价关系。那么关系 “R\{x, y\} 以及 S\{z, t\}''在积集 E × F 的元素 (x, z) 与 (y, t) 之间，关系是 E × F 上的一个等价关系，称为 R 与 S 的积，记作 R × S。关于 R × S 的每个等价类都是关于 R 的一个等价类与关于 S 的一个等价类的积。若 u 表示 E/R 的一般元素，v 表示 F/S 的一般元素，则 (u, v) → u × v 是从 (E/R) × (F/S) 到 (E × F)/(R × S) 的一个双射（称为典范双射）。
\end{enumerate}

